% ============================================================================
%  Tables for the coupled SWAT+ <-> MODFLOW 6 PFAS manuscript (Rogue River, MI)
%  Requires: \usepackage{booktabs}  (no siunitx; plain formatting for portability)
%  Every numeric value is grounded in the model artifacts / results.tex.
% ============================================================================

% ---------------------------------------------------------------------------
% Table 1 -- Study area and model configuration
% ---------------------------------------------------------------------------
\begin{table}[t]
  \centering
  \footnotesize
  \caption{Study area and coupled-model configuration. The surface-water
  (SWAT\textsuperscript{+}) and groundwater (MODFLOW~6) models are generated
  automatically from public data on the Rogue River watershed, Michigan.}
  \label{tab:studyarea}
  \begin{tabular}{@{}ll@{}}
    \toprule
    Quantity & Value \\
    \midrule
    \multicolumn{2}{@{}l}{\emph{Watershed and observations}}\\
    Watershed                         & Rogue River, Michigan \\
    USGS streamflow gauge             & 04118500 \\
    PFAS contamination source         & Wolverine World Wide / House Street site \\
    Source-zone anchor (measured GW PFOS) & cells $>10^{4}$~ng~L\textsuperscript{-1}; peak $1.5\times10^{6}$ \\
    Groundwater head observations     & 5{,}383 per-cell static water-table heads \\
    Groundwater PFOS observations     & 846 (downgradient + background) \\
    \addlinespace
    \multicolumn{2}{@{}l}{\emph{Surface-water model (SWAT\textsuperscript{+})}}\\
    Hydrography source                & NHDPlus High Resolution \\
    HRUs                              & 17{,}771 \\
    Channels (SWAT\textsuperscript{+} graph) & 590 \\
    \addlinespace
    \multicolumn{2}{@{}l}{\emph{Groundwater model (MODFLOW~6)}}\\
    Active cells                       & $\approx\!24{,}900$ \\
    Layers                             & 3 \\
    Cell size (uniform grid)           & 250~m \\
    Grid P\'eclet number               & $\sim$25 ($\Delta x/\alpha_L = 250/10$) \\
    SFR reaches                        & 1{,}506 \\
    SFR channels (order $\geq 2$)      & 273 \\
    Porosity                           & 0.30 \\
    Bulk density                       & 1800~kg~m\textsuperscript{-3} \\
    \addlinespace
    \multicolumn{2}{@{}l}{\emph{Transport simulation}}\\
    Flow field                         & calibrated steady state \\
    PFAS transport horizon             & 40-yr transient (adaptive time stepping) \\
    \bottomrule
  \end{tabular}
\end{table}

% ---------------------------------------------------------------------------
% Table 2 -- Calibration parameters (priors and posteriors)
% ---------------------------------------------------------------------------
\begin{table}[t]
  \centering
  \footnotesize
  \caption{Calibration parameters. Flow parameters are \emph{calibrated} by the
  PEST\texttt{++} iterative ensemble smoother (\textsc{ies}) against 5{,}383
  water-table heads with a baseflow constraint; values are the deployed
  posterior realization's global $\log_{10}$ multipliers on the prior field.
  Transport parameters are \emph{prescribed} from the PFAS literature, and the
  groundwater source is anchored to measurements rather than fitted.}
  \label{tab:calparams}
  \begin{tabular}{@{}llcc@{}}
    \toprule
    Parameter & Symbol / units & Prior range & Posterior / used value \\
    \midrule
    \multicolumn{4}{@{}l}{\emph{Flow (calibrated, PEST\texttt{++} \textsc{ies})}}\\
    Recharge multiplier            & $\log_{10}$ mult.            & $[-0.30,0.48]$  & $+0.188$ \\
    Vertical conductivity (anisotropy) & $k_v$, $\log_{10}$ mult.  & $[-0.8,0.8]$  & $+0.023$ \\
    GHB (divide) conductance       & $\log_{10}$ mult.            & $[-2.0,0.5]$  & $-1.169$ \\
    Drain conductance              & $\log_{10}$ mult.            & $[-2.0,0.5]$  & $+0.500$\textsuperscript{a} \\
    Well pumping / withdrawal       & $\log_{10}$ mult.            & $[-1.70,0.0]$  & $-1.264$ \\
    SFR streambed conductance      & $\log_{10}$ mult.           & $[-1.5,1.0]$  & $+0.693$ \\
    Hydraulic conductivity field    & 243 pilot points, $\log_{10}$ mult. & $[-0.8,0.8]$ & spatially variable \\
    \addlinespace
    \multicolumn{4}{@{}l}{\emph{Transport (prescribed, literature)}}\\
    Freundlich sorption coefficient & $K_f$                       & ---\textsuperscript{b} & 0.005 \\
    Freundlich exponent             & $n$                         & ---\textsuperscript{b} & 0.8 \\
    Longitudinal dispersivity       & $\alpha_L$ (m)              & ---\textsuperscript{b} & 10 \\
    Transverse dispersivity         & $\alpha_T$ (m)              & ---\textsuperscript{b} & 1 \\
    Molecular diffusion             & $D$ (m\textsuperscript{2}~s\textsuperscript{-1}) & ---\textsuperscript{b} & $1\times10^{-10}$ \\
    Source cap                      & ng~L\textsuperscript{-1}    & ---\textsuperscript{b} & $1\times10^{5}$ \\
    Source threshold                & ng~L\textsuperscript{-1}    & ---\textsuperscript{b} & $1\times10^{4}$ \\
    Ambient background              & ng~L\textsuperscript{-1}    & ---\textsuperscript{b} & 10 \\
    \bottomrule
  \end{tabular}

  \vspace{2pt}
  \begin{minipage}{\linewidth}\footnotesize
  \textsuperscript{a}The drain-conductance multiplier reaches its upper prior
  bound ($+0.5$), indicating the calibration favours stronger drainage than the
  prior centre; the remaining flow multipliers are interior.\quad
  \textsuperscript{b}Prescribed PFAS-transport literature value; not calibrated,
  so no prior range applies.
  \end{minipage}
\end{table}

% ---------------------------------------------------------------------------
% Table 3 -- Validation metrics (consolidated)
% ---------------------------------------------------------------------------
\begin{table}[t]
  \centering
  \footnotesize
  \caption{Consolidated validation metrics for the coupled system across the
  groundwater-flow, groundwater-transport, surface-water-transport, and joint
  compartments. Concentration skill is reported in $\log_{10}$ space (dex), the
  appropriate metric for values spanning orders of magnitude.}
  \label{tab:validation}
  \begin{tabularx}{\linewidth}{@{}>{\raggedright\arraybackslash}p{3.0cm}llcX@{}}
    \toprule
    Compartment / quantity & Metric & Value & $n$ & Method / notes \\
    \midrule
    Groundwater heads
      & Nash--Sutcliffe / RMSE & 0.91 / 5.6~m & 5{,}383
      & PEST\texttt{++} \textsc{ies}, baseflow-constrained \\
    Baseflow
      & gaining discharge & $+5.46$ vs 5.56 & ---
      & m\textsuperscript{3}~s\textsuperscript{-1}; observed gaining baseflow \\
    Groundwater PFOS (all predicted cells)
      & log-RMSE / within $\times10$ & 1.1~dex / 71\% & 65
      & 846 obs $\to$ 73 cells; 8 prescribed, 65 predicted \\
    Groundwater PFOS (above background)
      & log-RMSE / within $\times10$ & 1.4~dex / 46\% & 26
      & observed $>50$~ng~L\textsuperscript{-1}; 12/26 within $\times10$ \\
    Surface-water PFOS (source-bearing mainstem)
      & log-RMSE & 0.15~dex & 7
      & calibrated mainstem reaches \\
    Surface-water PFOS (all gauged reaches)
      & log-RMSE & 0.74~dex & 20
      & negative log Nash--Sutcliffe; tributary heterogeneity \\
    Joint SW+GW mainstem fit
      & log-RMSE; held-out RMSE & 0.07~dex; $6.2\!\to\!4.5$ & 7
      & NNLS $(L,g)$; held-out RMSE in ng~L\textsuperscript{-1} (Fig.~S14 reports MAE) \\
    \bottomrule
  \end{tabularx}
\end{table}

% ---------------------------------------------------------------------------
% Table 4 -- Joint SW + GW calibration, per mainstem reach
% ---------------------------------------------------------------------------
\begin{table}[t]
  \centering
  \footnotesize
  \caption{Joint surface-water + groundwater calibration along the Rogue
  mainstem (upstream $\to$ downstream). In-stream PFOS is linear in both the
  surface soil-loading $L$ and the groundwater source effectiveness $g$;
  fitting $(L,g)$ jointly by non-negative least squares gives $L=0.077$ (from a
  surface-only $L_0=0.11$) and $g=0.061$ (95\% CI $[0.022,0.100]$). The added
  groundwater term is significant against the surface-only ($g=0$) model
  ($F_{1,5}=16.2$, $p=0.01$) and survives leave-one-out cross-validation, and
  the groundwater fraction rises from $\sim$13\% in the headwaters to $\sim$53\%
  at the lower mainstem near the Wolverine source. Concentrations in
  ng~L\textsuperscript{-1}.}
  \label{tab:jointreach}
  \begin{tabular}{@{}rrrrrrr@{}}
    \toprule
    Channel & Observed & SW-only & Joint & SW part & GW part & GW frac.\ (\%) \\
    \midrule
    26 &  6.3 &  6.2 &  5.0 &  4.3 &  0.6 & 13 \\
    18 &  6.8 &  7.7 &  6.1 &  5.4 &  0.7 & 11 \\
    15 &  7.2 &  9.2 &  7.2 &  6.4 &  0.7 & 10 \\
    11 &  6.7 & 11.2 &  8.6 &  7.8 &  0.8 &  9 \\
    10 &  9.6 & 12.4 &  9.5 &  8.7 &  0.8 &  8 \\
     1 & 19.0 & 14.7 & 23.0 & 10.3 & 12.7 & 55 \\
     2 & 27.0 & 15.4 & 22.8 & 10.8 & 12.0 & 53 \\
    \bottomrule
  \end{tabular}
\end{table}
